\section{Bối cảnh}

Hoạt động bán lẻ hiện đại ngày càng được triển khai trên nhiều kênh, bao gồm điểm bán hàng
trực tiếp (POS) và các nền tảng thương mại điện tử như Shopee, TikTok Shop và Lazada. Khi
các kênh được quản lý độc lập, thông tin tồn kho dễ bị sai lệch và nhiều giao dịch có thể
cùng cạnh tranh trên một lượng hàng tồn thực tế.

Đề tài tập trung xây dựng một hệ thống quản lý bán hàng và tồn kho đa kênh, cung cấp góc
nhìn tập trung đối với sản phẩm, tồn kho, đơn hàng, giá vốn và hiệu quả bán hàng. Hệ thống
được định hướng theo mô hình phần mềm cung cấp dưới dạng dịch vụ (SaaS) đa khách hàng,
trong đó dữ liệu của từng tổ chức được cô lập về mặt logic.

\section{Phát biểu bài toán}

Các vấn đề chính cần giải quyết gồm:

\begin{itemize}
    \item bán vượt số lượng tồn khi nhiều kênh cùng bán một lượng hàng còn lại;
    \item số dư tồn kho không chính xác do điều chỉnh thủ công và dữ liệu giữa các kênh bị tách rời;
    \item khó tính giá vốn hàng bán (COGS) một cách nhất quán;
    \item đồng bộ đơn hàng trực tuyến với tồn kho tại cửa hàng bị chậm;
    \item thiếu khả năng theo dõi tập trung doanh thu, biến động tồn kho, lợi nhuận gộp và tốc độ bán.
\end{itemize}

Một giải pháp tin cậy phải phối hợp được việc tiếp nhận đơn hàng, giữ hàng, trừ tồn,
tính giá vốn và lập báo cáo trong điều kiện có nhiều giao dịch đồng thời.

\section{Mục tiêu đề tài}

Mục tiêu của hệ thống OISM gồm:

\begin{enumerate}
    \item quản lý tập trung sản phẩm, SKU, giá bán, tồn kho và đơn hàng;
    \item hỗ trợ xác thực đa khách hàng và phân quyền theo vai trò;
    \item duy trì sổ cái tồn kho chỉ ghi thêm để làm nguồn dữ liệu có khả năng truy vết;
    \item tính giá vốn bình quân gia quyền và lưu lại giá vốn tại thời điểm hoàn tất giao dịch;
    \item ngăn bán vượt bằng cơ chế giữ hàng trong giao dịch và khóa cơ sở dữ liệu;
    \item cung cấp quy trình bán hàng POS nhanh dưới dạng ứng dụng web tiến bộ (PWA);
    \item tiếp nhận đơn hàng trực tuyến, xử lý tác vụ nền và gửi thông báo thời gian thực;
    \item cung cấp báo cáo doanh thu, COGS, lợi nhuận gộp, tồn kho, tốc độ bán và cảnh báo tồn;
    \item hỗ trợ triển khai có thể tái lập bằng Docker và quy trình CI/CD.
\end{enumerate}

\section{Phạm vi}

Phạm vi hệ thống bao gồm xác thực, quản lý khách hàng và chi nhánh, danh mục sản phẩm,
nhóm sản phẩm, thương hiệu, biến thể, SKU, mã vạch, giá bán, giao dịch tồn kho, phiếu nhập,
chuyển kho, kiểm kê, tính giá vốn bình quân gia quyền, quản lý đơn hàng tập trung, giữ hàng,
thanh toán POS, báo cáo, mô phỏng webhook, thông báo SignalR và xử lý tác vụ nền bằng Hangfire.

Đề tài không xây dựng hạ tầng nội bộ của các sàn thương mại điện tử bên ngoài. Thay vào đó,
hành vi của kênh bên ngoài được biểu diễn bằng mô hình đơn hàng chuẩn và bộ mô phỏng webhook
phù hợp cho trình diễn và kiểm thử.

\section{Giải pháp đề xuất}

Giải pháp kết hợp các cơ chế kiến trúc và nghiệp vụ sau:

\begin{itemize}
    \item sử dụng Kiến trúc Sạch (Clean Architecture) để tách miền nghiệp vụ, ứng dụng, hạ tầng và giao diện;
    \item xây dựng nền tảng SaaS đa khách hàng với truy cập dữ liệu có nhận biết khách hàng;
    \item sử dụng Sổ cái tồn kho (Inventory Ledger) chỉ ghi thêm để truy vết mọi biến động;
    \item xây dựng Trung tâm tiếp nhận đơn hàng để chuẩn hóa sự kiện từ nhiều kênh;
    \item sử dụng máy trạng thái đơn hàng: \textit{Draft} $\rightarrow$ \textit{Reserved}
          $\rightarrow$ \textit{Confirmed} $\rightarrow$ \textit{Completed}, kèm các nhánh hủy;
    \item xây dựng bộ tính giá vốn bình quân gia quyền phục vụ định giá tồn kho và COGS;
    \item sử dụng cơ chế chống bán vượt dựa trên khóa dữ liệu khi giữ hàng đồng thời;
    \item xây dựng POS PWA cho thao tác bán hàng tại cửa hàng;
    \item sử dụng SignalR cho cập nhật thời gian thực và Hangfire cho tác vụ nền.
\end{itemize}

\section{Kết quả mong đợi}

Kết quả cuối cùng là một hệ thống có thể triển khai để trình diễn, gồm Web API .NET 8,
giao diện quản trị, POS PWA, tài liệu kỹ thuật, đặc tả API, Dockerfile, cấu hình Docker Compose
và quy trình CI/CD bằng GitHub Actions.

Đề tài đồng thời cung cấp bộ yêu cầu có khả năng truy vết, tài liệu kiến trúc, thiết kế chi tiết,
sơ đồ thiết kế, bằng chứng kiểm thử và hướng dẫn triển khai.
